\begin{abstract}
\emph{Objective:} To estimate and forecast subjective sleepiness on an ordinary laptop by combining a biomathematical model, occasional self-ratings and passive signals, and to test the approach on public data.
\emph{Methods:} BlinkTwin runs in a web browser; video never leaves the page. A Kalman filter tracks the person's offset from the Three-Process Model of alertness, updated by Karolinska Sleepiness Scale (KSS) check-ins. A 3-min psychomotor vigilance test, blink duration and PERCLOS enter through forward models fitted within person on DROZY (14 adults), each with its own personal anchor state. We replayed the released model on DROZY leaving one subject out, ran its webcam pipeline on DROZY video against the electro-oculogram (EOG), analysed OpenNeuro ds004902 (33 adults), and ran a simulation with misspecified scenarios.
\emph{Results:} On DROZY, the mean absolute error of predicted KSS was 1.73 points for the prior, 1.51 with one check-in, 1.47 with passive signals and 1.38 with all three; intervals overlapped and a protocol-aware baseline gave 1.39. Webcam blink rates agreed with EOG (Spearman $\rho = 0.71$) and blink duration tracked KSS within person ($r = 0.58$), but PERCLOS failed on near-infrared video. On ds004902, offsets persisted weakly ($r = 0.29$) and a field-refitted KSS mapping lowered the prior's error from 1.69 to 1.46. In simulation, anchored fusion lowered the error of check-ins alone by 0.03 points; fusion without anchors failed.
\emph{Conclusion:} Check-ins set the level; passive signals can be added safely only with their own anchors, and their gains were small.
\emph{Significance:} An open, on-device implementation and a pre-specified protocol for testing personal sleepiness models in daily life.
\end{abstract}

\begin{IEEEkeywords}
Digital twin, Kalman filter, Karolinska Sleepiness Scale, PERCLOS, psychomotor vigilance test, sleepiness, Three-Process Model, webcam.
\end{IEEEkeywords}
